% GENERATED FILE. Edit canonical lessons, then run npm run generate:books.
% canonical-source-hash: fnv1a64:9207a732fa14b6ca
% canonical-lessons: JA-C459-daigaku, JA-C459-ryuugaku-suru, JA-C459-kenkyuu-suru, JA-C459-senmon, JA-C459-shinro

\chapter{Studying Further}
\label{ch:ja-studying-further}
\hlchaptermodality{459}

\begin{chapteropening}
\textbf{By the end of this chapter:} \emph{I can talk about university and research.}

Complete the last lesson of chapter 459: I can talk about university and research.
\end{chapteropening}

\section[\texorpdfstring{\ja{だいがく} (daigaku)}{だいがく (daigaku)}]{\ja{だいがく} (daigaku) — a university}

\label{lesson:JA-C459-daigaku}



\begin{quote}
Before the new one: say the Japanese for a novelist, then the Japanese for a craftsman.
\end{quote}

\subsection*{You'll want to know: \ja{だいがく}}
\textbf{\ja{だいがく}} — \emph{daigaku} — \textquotedblleft{}a university\textquotedblright{}.

\begin{quote}
\textbf{\ja{あには} \ja{だいがくで} \ja{べんきょうしています。}} — \emph{ani wa daigaku de benkyō shite imasu} — \textquotedblleft{}My elder brother studies at university.\textquotedblright{}
\end{quote}

\begin{cousinweb}
\textbf{\ja{だいがく}} (\emph{daigaku}) is Sino-Japanese (kango): two Chinese characters read with their on'yomi, readings borrowed from Chinese. They are \emph{dai}, \textquotedblleft{}large\textquotedblright{} (Middle Chinese \emph{dajH}) and \emph{gaku}, \textquotedblleft{}study\textquotedblright{} (\emph{haewk}) (KANJIDIC2; Middle Chinese in Baxter's transcription).
\end{cousinweb}

\begin{grammarlens}[title={What you've built}]
One more word for describing what you lived and felt.
\end{grammarlens}

\subsection*{Your turn}
\begin{itemize}
\raggedright
  \item \emph{Say it:} \emph{daigaku}
  \item \emph{Say it:} \emph{daigaku}, once more
  \item \emph{Say it:} \emph{ani wa daigaku de benkyō shite imasu} — My elder brother studies at university.
  \item \emph{Recall:} say the Japanese for a novelist, then the Japanese for a craftsman, then say \emph{daigaku} again
\end{itemize}

\begin{tcolorbox}[breakable,colback=teal!4,colframe=teal!35!black,title={Before you move on}]
What does \ja{だいがく} mean? (\textquotedblleft{}A university\textquotedblright{}.) Say it once more.
\end{tcolorbox}
\section[\texorpdfstring{\ja{りゅうがくする} (ryūgaku suru)}{りゅうがくする (ryūgaku suru)}]{\ja{りゅうがくする} (ryūgaku suru) — to study abroad}

\label{lesson:JA-C459-ryuugaku-suru}



\begin{quote}
Before the new one: say the Japanese for a craftsman, then the Japanese for a university.
\end{quote}

\subsection*{You'll want to know: \ja{りゅうがくする}}
\textbf{\ja{りゅうがくする}} — \emph{ryūgaku suru} — \textquotedblleft{}to study abroad\textquotedblright{}.

\begin{quote}
\textbf{\ja{ともだちは} \ja{らいねん} \ja{りゅうがくします。}} — \emph{tomodachi wa rainen ryūgaku shimasu} — \textquotedblleft{}My friend is studying abroad next year.\textquotedblright{}
\end{quote}

Polite \textbf{\ja{りゅうがくします}}, \emph{ryūgakushimasu}; a \ja{する} verb. Past (polite): \textbf{\ja{あねが} \ja{りゅうがくしました。}}, \emph{ane ga ryūgaku shimashita}.

\begin{cousinweb}
\textbf{\ja{りゅうがくする}} (\emph{ryūgaku suru}) is Sino-Japanese (kango) \emph{ryūgaku} with the native verb \emph{suru}, \textquotedblleft{}to do\textquotedblright{}: two Chinese characters read with their on'yomi, readings borrowed from Chinese. They are \emph{ryū}, \textquotedblleft{}detain\textquotedblright{} (Middle Chinese \emph{ljuw}) and \emph{gaku}, \textquotedblleft{}study\textquotedblright{} (\emph{haewk}) (KANJIDIC2; Middle Chinese in Baxter's transcription).
\end{cousinweb}

\begin{grammarlens}[title={What you've built}]
One more word for describing what you lived and felt.
\end{grammarlens}

\subsection*{Your turn}
\begin{itemize}
\raggedright
  \item \emph{Say it:} \emph{ryūgaku suru}
  \item \emph{Say it:} \emph{ryūgaku suru}, once more
  \item \emph{Say it:} \emph{tomodachi wa rainen ryūgaku shimasu} — My friend is studying abroad next year.
  \item \emph{Recall:} say the Japanese for a craftsman, then the Japanese for a university, then say \emph{ryūgaku suru} again
\end{itemize}

\begin{tcolorbox}[breakable,colback=teal!4,colframe=teal!35!black,title={Before you move on}]
What does \ja{りゅうがくする} mean? (\textquotedblleft{}To study abroad\textquotedblright{}.) Say it once more.
\end{tcolorbox}
\section[\texorpdfstring{\ja{けんきゅうする} (kenkyū suru)}{けんきゅうする (kenkyū suru)}]{\ja{けんきゅうする} (kenkyū suru) — to do research}

\label{lesson:JA-C459-kenkyuu-suru}



\begin{quote}
Before the new one: say the Japanese for a university, then the Japanese for to study abroad.
\end{quote}

\subsection*{You'll want to know: \ja{けんきゅうする}}
\textbf{\ja{けんきゅうする}} — \emph{kenkyū suru} — \textquotedblleft{}to do research\textquotedblright{}.

\begin{quote}
\textbf{\ja{かれは} \ja{だいがくで} \ja{しょくぶつを} \ja{けんきゅうしています。}} — \emph{kare wa daigaku de shokubutsu o kenkyū shite imasu} — \textquotedblleft{}He does research on plants at a university.\textquotedblright{}
\end{quote}

Polite \textbf{\ja{けんきゅうします}}, \emph{kenkyūshimasu}; a \ja{する} verb. Past (polite): \textbf{\ja{せんせいが} \ja{どうぶつを} \ja{けんきゅうしました。}}, \emph{sensei ga dōbutsu o kenkyū shimashita}.

\begin{cousinweb}
\textbf{\ja{けんきゅうする}} (\emph{kenkyū suru}) is Sino-Japanese (kango) \emph{kenkyū} with the native verb \emph{suru}, \textquotedblleft{}to do\textquotedblright{}: two Chinese characters read with their on'yomi, readings borrowed from Chinese. They are \emph{ken}, \textquotedblleft{}polish\textquotedblright{} (Middle Chinese \emph{ngen}) and \emph{kyū}, \textquotedblleft{}research\textquotedblright{} (\emph{kjuwH}) (KANJIDIC2; Middle Chinese in Baxter's transcription).
\end{cousinweb}

\begin{grammarlens}[title={What you've built}]
One more word for describing what you lived and felt.
\end{grammarlens}

\subsection*{Your turn}
\begin{itemize}
\raggedright
  \item \emph{Say it:} \emph{kenkyū suru}
  \item \emph{Say it:} \emph{kenkyū suru}, once more
  \item \emph{Say it:} \emph{kare wa daigaku de shokubutsu o kenkyū shite imasu} — He does research on plants at a university.
  \item \emph{Recall:} say the Japanese for a university, then the Japanese for to study abroad, then say \emph{kenkyū suru} again
\end{itemize}

\begin{tcolorbox}[breakable,colback=teal!4,colframe=teal!35!black,title={Before you move on}]
What does \ja{けんきゅうする} mean? (\textquotedblleft{}To do research\textquotedblright{}.) Say it once more.
\end{tcolorbox}
\section[\texorpdfstring{\ja{せんもん} (senmon)}{せんもん (senmon)}]{\ja{せんもん} (senmon) — one's field of study; a specialty}

\label{lesson:JA-C459-senmon}



\begin{quote}
Before the new one: say the Japanese for to study abroad, then the Japanese for to do research.
\end{quote}

\subsection*{You'll want to know: \ja{せんもん}}
\textbf{\ja{せんもん}} — \emph{senmon} — \textquotedblleft{}one's field of study; a specialty\textquotedblright{}.

\begin{quote}
\textbf{\ja{あねの} \ja{せんもんは} \ja{おんがくです。}} — \emph{ane no senmon wa ongaku desu} — \textquotedblleft{}My elder sister's field is music.\textquotedblright{}
\end{quote}

\begin{cousinweb}
\textbf{\ja{せんもん}} (\emph{senmon}) is Sino-Japanese (kango): two Chinese characters read with their on'yomi, readings borrowed from Chinese. They are \emph{sen}, \textquotedblleft{}specialty\textquotedblright{} (Middle Chinese \emph{tsywen}) and \emph{mon}, \textquotedblleft{}gate\textquotedblright{} (\emph{mwon}) (KANJIDIC2; Middle Chinese in Baxter's transcription).
\end{cousinweb}

\begin{grammarlens}[title={What you've built}]
One more word for describing what you lived and felt.
\end{grammarlens}

\subsection*{Your turn}
\begin{itemize}
\raggedright
  \item \emph{Say it:} \emph{senmon}
  \item \emph{Say it:} \emph{senmon}, once more
  \item \emph{Say it:} \emph{ane no senmon wa ongaku desu} — My elder sister's field is music.
  \item \emph{Recall:} say the Japanese for to study abroad, then the Japanese for to do research, then say \emph{senmon} again
\end{itemize}

\begin{tcolorbox}[breakable,colback=teal!4,colframe=teal!35!black,title={Before you move on}]
What does \ja{せんもん} mean? (\textquotedblleft{}One's field of study; a specialty\textquotedblright{}.) Say it once more.
\end{tcolorbox}
\section[\texorpdfstring{\ja{しんろ} (shinro)}{しんろ (shinro)}]{\ja{しんろ} (shinro) — a career path; one's course}

\label{lesson:JA-C459-shinro}



\begin{quote}
Before the new one: say the Japanese for to do research, then the Japanese for one's field of study.
\end{quote}

\subsection*{You'll want to know: \ja{しんろ}}
\textbf{\ja{しんろ}} — \emph{shinro} — \textquotedblleft{}a career path; one's course\textquotedblright{}.

\begin{quote}
\textbf{\ja{せんせいと} \ja{しんろを} \ja{はなしあいました。}} — \emph{sensei to shinro o hanashiaimashita} — \textquotedblleft{}I talked over my career path with my teacher.\textquotedblright{}
\end{quote}

\begin{cousinweb}
\textbf{\ja{しんろ}} (\emph{shinro}) is Sino-Japanese (kango): two Chinese characters read with their on'yomi, readings borrowed from Chinese. They are \emph{shin}, \textquotedblleft{}advance\textquotedblright{} (Middle Chinese \emph{tsinH}) and \emph{ro}, \textquotedblleft{}path\textquotedblright{} (\emph{luH}) (KANJIDIC2; Middle Chinese in Baxter's transcription).
\end{cousinweb}

\begin{grammarlens}[title={What you've built}]
One more word for describing what you lived and felt.
\end{grammarlens}

\subsection*{Your turn}
\begin{itemize}
\raggedright
  \item \emph{Say it:} \emph{shinro}
  \item \emph{Say it:} \emph{shinro}, once more
  \item \emph{Say it:} \emph{sensei to shinro o hanashiaimashita} — I talked over my career path with my teacher.
  \item \emph{Recall:} say the Japanese for to do research, then the Japanese for one's field of study, then say \emph{shinro} again
\end{itemize}

\begin{tcolorbox}[breakable,colback=teal!4,colframe=teal!35!black,title={Before you move on}]
What does \ja{しんろ} mean? (\textquotedblleft{}A career path; one's course\textquotedblright{}.) Say it once more.
\end{tcolorbox}
